\documentclass[12pt]{article}
\usepackage{setspace}
\onehalfspacing

\usepackage{iftex}
\ifPDFTeX
  \usepackage[utf8]{inputenc}
  \usepackage[T1]{fontenc}
\else
  \usepackage{fontspec}
\fi
\usepackage[spanish, es-tabla]{babel}
\usepackage[margin=1in]{geometry}
\usepackage{enumitem}
\usepackage{hyperref}
\usepackage{microtype}
\usepackage{biblatex}
\addbibresource{refs.bib}

\hypersetup{
  colorlinks=true,
  linkcolor=black,
  urlcolor=blue,
  citecolor=black
}

\setlist[itemize]{leftmargin=*, itemsep=0.35em}
\newcommand{\institucion}{Universidad de Costa Rica}
\newcommand{\campus}{Ciudad Universitaria Rodrigo Facio}
\newcommand{\facultad}{Facultad de Ciencias Básicas}
\newcommand{\escuela}{Escuela de Física}
\newcommand{\curso}{Mecánica Estadística Computacional}
\newcommand{\tituloanteproyecto}{Anteproyecto de Investigación}
\newcommand{\titulo}{Simulación Monte Carlo de Casos Nuevos de COVID-19 por Zona Geográfica}
\newcommand{\autor}{Gustavo A. Calvo Gutiérrez}
\newcommand{\lugarfecha}{San José, 2026.}

\begin{document}

\begin{titlepage}
  \begin{center}
    \LARGE \textbf{\textsc{\institucion}}
    \Large
    \\[1cm]
    \textsc{\campus}
    \\[0.1cm]
    \textsc{\facultad}
    \\[0.1cm]
    \textsc{\escuela} \\
    \vspace*{20 mm}
    \LARGE{\textsc{\textbf{\curso}}}
  \end{center}

  \vspace{2cm}

  \hrulefill
  \begin{center}
    \large \textbf{\textsc{\tituloanteproyecto}}
    \\[0.5cm]
    \textsc{\titulo}
  \end{center}
  \hrulefill

  \vfill

  \begin{center}
    \large \textbf{\autor}
  \end{center}

  \vspace{1cm}

  \begin{center}
    \textit{\lugarfecha}
  \end{center}
\end{titlepage}

\section{Pregunta de Investigación y Justificación}

\subsection{Pregunta de Investigación}

\noindent
\textit{
  ¿Cómo puede un modelo de simulación Monte Carlo, aplicado por zona
  geográfica y basado en el procedimiento propuesto por Xie \cite{Xie2020},
  estimar la cantidad esperada de casos nuevos de COVID-19 a partir de series
  temporales observadas?
}

\subsection{Justificación}

El estudio de la propagación de COVID-19 continúa siendo relevante como caso de
referencia para comprender fenómenos epidémicos donde la transmisión depende de
interacciones sociales y de variaciones temporales en la dinámica de contagio.
Los modelos estocásticos ofrecen una ventaja importante frente a aproximaciones
puramente deterministas, ya que permiten representar explícitamente la
incertidumbre asociada a la ocurrencia de nuevos contagios y a la variabilidad
temporal de la transmisión.

El trabajo de Xie \cite{Xie2020} plantea una estrategia basada en simulación
Monte Carlo para describir la dinámica de contagio mediante tasas de infección
variables en el tiempo, junto con distribuciones probabilísticas para modelar
tanto la cantidad de personas contagiadas por cada caso activo como el momento
en que dichos contagios ocurren.

El presente anteproyecto propone tomar ese artículo como punto de partida
metodológico para realizar simulaciones por zona geográfica y estimar, a partir
de múltiples corridas aleatorias, la cantidad esperada de casos nuevos en cada
periodo. La meta no es evaluar la efectividad de medidas sanitarias ni comparar
intervenciones, sino construir una herramienta de simulación capaz de generar
trayectorias plausibles de contagio y resumirlas mediante su valor esperado.
Este enfoque es pertinente porque permite vincular la simulación con series
epidemiológicas observadas y estudiar la variabilidad inherente al proceso de
transmisión. Para ello, el estudio puede apoyarse en fuentes comparables de
datos epidemiológicos, como \textit{Our World in Data}
\cite{owid-coronavirus}.

\newpage

\section{Objetivos}

\subsection{Objetivo General}

Desarrollar un modelo de simulación Monte Carlo, basado en el procedimiento
propuesto por Xie \cite{Xie2020}, para estimar la cantidad esperada de casos
nuevos de COVID-19 en distintas zonas geográficas.

\subsection{Objetivos Específicos}

\begin{itemize}
  \item Adaptar el procedimiento Monte Carlo de Xie \cite{Xie2020} a un
        conjunto de datos temporales de COVID-19 organizado por zona geográfica.
  \item Recopilar y estructurar series de tiempo con variables relevantes,
        tales como casos nuevos, personas infectadas, recuperadas, casos
        activos y población de referencia por zona, utilizando fuentes
        comparables \cite{owid-coronavirus}.
  \item Modelar la probabilidad diaria de contagio como un componente aleatorio
        dependiente del tiempo y calibrado con la evolución observada de los
        contagios en cada zona.
  \item Comparar las trayectorias simuladas con los datos reales observados
        para identificar discrepancias, patrones recurrentes y posibles ajustes
        del modelo.
  \item Estimar, a partir de múltiples simulaciones, el valor esperado de
        casos nuevos por periodo en cada zona analizada.
\end{itemize}

\newpage

\section{Marco Teórico}

La simulación de epidemias puede abordarse desde modelos deterministas, como
los de tipo SIR, o desde modelos estocásticos que incorporan explícitamente la
incertidumbre del proceso de contagio.
Los modelos estocásticos consideran que la transmisión
ocurre a través de eventos aleatorios cuya realización concreta puede variar
incluso bajo condiciones iniciales similares. Esta segunda perspectiva resulta
especialmente útil cuando se desea representar fluctuaciones diarias,
heterogeneidad espacial y cambios abruptos en el comportamiento social.

En el trabajo de Xie \cite{Xie2020},
la propuesta modela la propagación de COVID-19 como un proceso de puntos en el
que cada persona infecciosa puede generar un número aleatorio de nuevos
contagios. La cantidad de infecciones secundarias se representa mediante una
distribución de Poisson de una tasa de infección \(R_t\),
mientras que el momento en que ocurre cada nuevo contagio se modela mediante
una distribución binomial negativa con media \(\mu_T\). Esta construcción
permite simular tanto la cantidad como el tiempo de aparición de los nuevos
casos.

Un punto central del modelo es que \(R_t\) no se considera constante. La tasa
de transmisión cambia con el tiempo para reflejar variaciones en el contacto
social y el comportamiento de la población. En el artículo de referencia, la
selección de patrones para \(R_t\) se apoya en el ajuste a datos observados
\cite{Xie2020}. Esta idea es especialmente valiosa para el presente proyecto,
porque permite representar cambios en la intensidad de la transmisión sin
abandonar el carácter probabilístico del modelo.

Desde el punto de vista empírico, la simulación requiere series de tiempo
comparables. No basta con observar el número acumulado de casos; es necesario
considerar variables que cambian día con día, tales como el número de personas
activamente infectadas, recuperadas, la población potencialmente expuesta y los
casos nuevos reportados en cada zona. Para documentar estos factores, resultan
particularmente útiles bases de datos como \textit{Our World in Data}
\cite{owid-coronavirus}, que reúne indicadores epidemiológicos comparables.

El proyecto también se apoya en la idea de validación por contraste entre
simulación y observación.
Si el modelo captura adecuadamente la tendencia general de contagios y la
variación de los casos nuevos en cada zona, entonces puede utilizarse para
obtener estimaciones esperadas a partir de muchas realizaciones del proceso
aleatorio. Si no lo hace, las discrepancias mismas aportan información sobre
variables omitidas o supuestos que deben refinarse.

\newpage

\section{Aplicación de la Técnica de Monte Carlo}

\subsection{Recopilación y Organización de Datos}

La primera etapa consistirá en construir una base de datos temporal para zonas
específicas de análisis. Se buscará registrar, para cada unidad espacial y cada
fecha, variables como casos nuevos, casos activos, personas recuperadas,
fallecimientos y población total o estimada. La organización en series de
tiempo permitirá observar la evolución del problema dentro de cada zona y usar
esa información como base para calibrar el modelo estocástico.

\subsection{Construcción del Modelo Estocástico}

Tomando como referencia el trabajo de Xie \cite{Xie2020}, se propondrá un
modelo en el que cada persona infectada tenga una probabilidad de contagiar a
otras personas en cada día del periodo de estudio. La parte aleatoria del
problema se representará precisamente en esa posibilidad diaria de transmisión.
En lugar de asumir que todos los contagios siguen una trayectoria fija, se
permitirá que el número de infecciones nuevas surja de distribuciones
probabilísticas calibradas con los datos observados.

De manera conceptual, el procedimiento será el siguiente:

\begin{itemize}
  \item Se parte de un número inicial de personas infectadas en una zona y
        fecha determinadas.
  \item Para cada día, se estima una tasa de transmisión asociada al contexto
        del periodo y calibrada con el comportamiento observado de los datos en
        esa zona.
  \item Cada caso activo puede generar nuevos contagios de forma aleatoria, y
        el momento exacto de esos contagios también se modela
        probabilísticamente.
  \item El proceso se repite durante toda la ventana temporal de estudio para
        obtener trayectorias simuladas de casos nuevos, casos activos y
        acumulados.
\end{itemize}

Una vez generadas múltiples trayectorias simuladas, estas se compararán con los
datos reales de cada zona para verificar si el modelo reproduce de forma
razonable la evolución observada. A partir del conjunto de resultados
obtenidos, se calculará el promedio de casos nuevos por periodo, el cual se
interpretará como una estimación del valor esperado buscado.

\newpage

\printbibliography

\end{document}
